% Mathematics Bootcamp --- One Statement, Three Proofs
% A single implication about parity, proved directly, by contrapositive,
% and by contradiction, with every step justified.
%
% Self-contained: this file needs no other file in the repository.
% Compile with pdflatex, twice, or with latexmk -pdf.

\documentclass[12pt]{article}

\usepackage[margin=1in]{geometry}

\usepackage{
  amsmath,
  amssymb,
  amsfonts,
  amsthm,
}

\usepackage[T1]{fontenc}
\usepackage{lmodern}

\usepackage{setspace}
\usepackage{float}
\usepackage{enumitem}
\usepackage{booktabs}

\onehalfspacing

% ---------------------------------------------------------------------------
% Same additions as the Day 1 exercise handout: a heading for each item, a
% "Need to Show" block, and one level of indentation for a structured proof.
% ---------------------------------------------------------------------------

% \method{number}{name}{one-line description}
\newcommand{\method}[3]{%
  \par\penalty-250\bigskip
  \noindent\textbf{\large Proof #1}\quad\textnormal{(#2)}%
  \par\nobreak\smallskip
  \noindent\textit{#3}%
  \par\nobreak\smallskip
  \noindent\rule{\linewidth}{0.4pt}%
  \par\nobreak\smallskip}

% What the proof is obliged to produce, written before any of it is written.
\newenvironment{nts}
  {\par\medskip\noindent\textbf{Need to Show:}\par\nobreak\smallskip}
  {\par\medskip}

% One level of indentation. Nest to indent further: the structure of the
% argument is carried by the indentation, as in the Day 1 slides.
\newlength{\plevel}\setlength{\plevel}{2.25em}
\newenvironment{pstep}
  {\begin{list}{}{%
     \setlength{\leftmargin}{\plevel}%
     \setlength{\rightmargin}{0pt}%
     \setlength{\labelwidth}{0pt}\setlength{\labelsep}{0pt}%
     \setlength{\itemindent}{0pt}\setlength{\listparindent}{0pt}%
     \setlength{\topsep}{0pt}\setlength{\partopsep}{0pt}%
     \setlength{\parsep}{0pt}\setlength{\itemsep}{0pt}}%
   \item\relax}
  {\end{list}}

% Right-flushed justification for a line of a proof. Ends the line.
\newcommand{\by}[1]{%
  \unskip\nobreak\hfil\penalty50\hskip1em\null\nobreak\hfil
  \textnormal{[#1]}%
  {\parfillskip=0pt \finalhyphendemerits=0 \par}}

\newenvironment{remark}
  {\par\medskip\noindent\textit{Remark.}\enspace}
  {\par\medskip}

% Contradiction. The single arrow is the one the Day 1 slides use for
% implication, so this reads as "then, and then back the other way".
\newcommand{\contra}{\ensuremath{\rightarrow\leftarrow}}

\title{\textbf{Mathematics Bootcamp --- One Statement, Three Proofs}\\[.4em]
       \large Direct, contrapositive, and contradiction on a single claim}
\author{Miguel V\'azquez-Carrero}
\date{Companion to Day 1: Formal Language and Proof Structure\\[.3em]
      \normalsize Compilation Date: \today}

\begin{document}

\maketitle

\noindent
One claim, proved three times. The arithmetic is deliberately trivial, so that
nothing on the page competes with the thing being taught: what changes between
the three methods is not the mathematics but the \emph{bookkeeping} --- what you
are allowed to assume, and what you are obliged to produce.

\section*{The claim}

\begin{quote}
\textbf{Claim.} Let $n\in\mathbb{Z}$. If $n$ is even, then $n+3$ is odd.
\end{quote}

\noindent
Write
\[
P:\ \text{``$n$ is even''},\qquad Q:\ \text{``$n+3$ is odd''},
\]
so the claim is $P\rightarrow Q$, to be proved for every $n\in\mathbb{Z}$.

\begin{quote}
\textbf{Definition.}
An integer $n$ is \emph{even} if $n=2k$ for some $k\in\mathbb{Z}$, and
\emph{odd} if $n=2k+1$ for some $k\in\mathbb{Z}$.
\end{quote}

\noindent
Read the definition as a contract in both directions. To \emph{use} ``$n$ is
even'' is to obtain a $k$; to \emph{prove} ``$n$ is even'' is to exhibit one.
Every step below is one of those two moves, an application of a background fact,
or ordinary algebra.

The only arithmetic assumed: $\mathbb{Z}$ is closed under addition, subtraction
and multiplication.

\section*{Two background facts}

The three proofs do not all rest on the same ground, and it is worth knowing
which fact each one spends. Both facts below are about parity alone; neither
mentions the claim, and both are taken as given here.

\begin{quote}
\textbf{Fact 1 (every integer has a parity).}
Every $n\in\mathbb{Z}$ is even or odd.

\medskip
\textbf{Fact 2 (no integer has two parities).}
No $n\in\mathbb{Z}$ is both even and odd.
\end{quote}

\noindent
Neither is proved here. They are the case $d=2$ of the statement that dividing
by $d$ leaves exactly one remainder in $\{0,\dots,d-1\}$, which is where they
belong; assume them today and take them up when the tools are available. Fact~1
in particular is the sort of ``obviously true'' statement whose proof needs
induction, which we have not reached yet --- a good illustration that obvious
and easy to prove are different properties.

What matters here is only how they get \emph{used}:
\begin{enumerate}[label=(\arabic*),leftmargin=2.4em,topsep=.3em,itemsep=.2em]
  \item Fact 1 turns ``$n$ is not odd'' into ``$n$ is even'' --- a negative
        hypothesis, which offers nothing to compute with, into a positive one
        that hands you a $k$;
  \item Fact 2 turns ``$n$ is odd'' into ``$n$ is not even'' --- letting a
        negative conclusion be reached by proving a positive statement.
\end{enumerate}
Between them, ``not odd'' and ``even'' are interchangeable, and so are ``not
even'' and ``odd''.

\begin{remark}
Keep an eye on where these get used below. Fact~1 appears twice --- once in
each of Proofs 2 and 3, in both cases on the very first line --- and never in
Proof~1. That is not an accident, and it is the point the three proofs together
are making.
\end{remark}

\section*{Why there are three methods at all}

Each method is licensed by a logical equivalence, and each equivalence is a
column of a truth table that the Day 1 slides already carry:
\[
\begin{array}{cc@{\qquad}cc@{\qquad}ccc}
\toprule
P & Q & \lnot P & \lnot Q &
P\rightarrow Q & \lnot Q\rightarrow\lnot P & \lnot(P\land\lnot Q)\\
\midrule
T & T & F & F & T & T & T\\
T & F & F & T & F & F & F\\
F & T & T & F & T & T & T\\
F & F & T & T & T & T & T\\
\bottomrule
\end{array}
\]
The last three columns are identical. So proving any one of them proves the
other two, and you may pick whichever is easiest to attack:
\begin{enumerate}[label=(\arabic*),leftmargin=2.4em,topsep=.3em,itemsep=.2em]
  \item $P\rightarrow Q$ --- assume $P$, produce $Q$;
  \item $\lnot Q\rightarrow\lnot P$ --- assume $\lnot Q$, produce $\lnot P$;
  \item $\lnot(P\land\lnot Q)$ --- assume both $P$ and $\lnot Q$, produce a
        contradiction.
\end{enumerate}
Nothing distinguishes them logically. They differ only in what lands in your
hands when you start, and that is a practical matter, not a matter of truth.

\section*{The three skeletons}

Written as programs, before any of the bodies are known. The indentation is not
decoration: a line is available exactly to the lines indented under it, and a
block is closed by the line that returns to the outer margin. The bracket is the
hole you still owe.

\begin{quote}
\textbf{Direct}\par
Let $n\in\mathbb{Z}$ be arbitrary.\par
\hspace*{2.25em}Suppose $n$ is even.\par
\hspace*{4.5em}[Proof that $n+3$ is odd goes here.]\par
\hspace*{2.25em}Thus, if $n$ is even then $n+3$ is odd.\par
Thus, for all $n\in\mathbb{Z}$, if $n$ is even then $n+3$ is odd.
\end{quote}

\begin{quote}
\textbf{Contrapositive}\par
Let $n\in\mathbb{Z}$ be arbitrary.\par
\hspace*{2.25em}Suppose $n+3$ is not odd.\par
\hspace*{4.5em}[Proof that $n$ is not even goes here.]\par
\hspace*{2.25em}Thus, if $n+3$ is not odd then $n$ is not even.\par
Thus $\lnot Q\rightarrow\lnot P$ for all $n$, which is $P\rightarrow Q$.
\end{quote}

\begin{quote}
\textbf{Contradiction}\par
Let $n\in\mathbb{Z}$ be arbitrary.\par
\hspace*{2.25em}Suppose $n$ is even and $n+3$ is not odd.\par
\hspace*{4.5em}[Proof of a contradiction goes here.]\par
\hspace*{2.25em}Thus $n$ even and $n+3$ not odd is impossible.\par
Thus $\lnot(P\land\lnot Q)$ for all $n$, which is $P\rightarrow Q$.
\end{quote}

\noindent
Write these outer lines first. They are forced by the shape of the statement and
the method alone --- before you have any idea why the claim is true.

\newpage
\section*{The three proofs}

% ===========================================================================
\method{1}{direct}{Assume $P$; produce $Q$.}

\begin{nts}
Assuming $n$ is even, produce an integer $j$ with $n+3=2j+1$. That is the whole
obligation: ``$n+3$ is odd'' is by definition the existence of such a $j$, so
the proof is finished exactly when one has been exhibited.
\end{nts}

\begin{proof}[Proof 1]
Let $n\in\mathbb{Z}$ be arbitrary.
\begin{pstep}
Suppose $n$ is even.
\begin{pstep}
$n=2k$ for some $k\in\mathbb{Z}$. \by{definition of even}

$n+3=2k+3$. \by{add $3$ to both sides}

$2k+3=2k+2+1=2(k+1)+1$. \by{algebra}

$k+1\in\mathbb{Z}$. \by{closure under addition}

So $n+3=2j+1$ with $j=k+1\in\mathbb{Z}$: $n+3$ is odd.
\by{definition of odd}
\end{pstep}
Thus, if $n$ is even then $n+3$ is odd.
\end{pstep}
Thus, for all $n\in\mathbb{Z}$, if $n$ is even then $n+3$ is odd.
\end{proof}

\begin{remark}
The witness is $j=k+1$, and producing it is the entire proof. Notice what was
\emph{not} used: neither fact. This argument would survive in a world where
integers could be both even and odd, or neither. It is the cheapest of the
three, and on a statement this shape it is the one to reach for.
\end{remark}

% ===========================================================================
\method{2}{contrapositive}{Assume $\lnot Q$; produce $\lnot P$.}

\begin{nts}
Prove $\lnot Q\rightarrow\lnot P$: if $n+3$ is not odd, then $n$ is not even.
Both ends need unfolding before there is anything to do:
\begin{enumerate}[label=(\alph*),leftmargin=2.4em,topsep=.2em,itemsep=.1em]
  \item the hypothesis ``$n+3$ is not odd'' is unusable as it stands --- it
        offers no $k$ --- and Fact~1 converts it into ``$n+3$ is even'', which
        does;
  \item the conclusion ``$n$ is not even'' is a negative, and Fact~2 lets it be
        reached by proving the positive statement ``$n$ is odd'' instead.
\end{enumerate}
So the work is: assuming $n+3=2m$, produce an integer $j$ with $n=2j+1$.
\end{nts}

\begin{proof}[Proof 2]
Let $n\in\mathbb{Z}$ be arbitrary.
\begin{pstep}
Suppose $n+3$ is not odd.
\begin{pstep}
$n+3$ is even. \by{Fact 1: not odd, so even}

$n+3=2m$ for some $m\in\mathbb{Z}$. \by{definition of even}

$n=2m-3$. \by{subtract $3$ from both sides}

$2m-3=2m-4+1=2(m-2)+1$. \by{algebra}

$m-2\in\mathbb{Z}$. \by{closure under subtraction}

So $n=2j+1$ with $j=m-2\in\mathbb{Z}$: $n$ is odd.
\by{definition of odd}

$n$ is not even. \by{Fact 2: odd, so not even}
\end{pstep}
Thus, if $n+3$ is not odd then $n$ is not even.
\end{pstep}
Thus $\lnot Q\rightarrow\lnot P$ holds for every $n\in\mathbb{Z}$; equivalently,
for every $n\in\mathbb{Z}$, if $n$ is even then $n+3$ is odd.
\by{contraposition}
\end{proof}

\begin{remark}
The last line is not a flourish --- it is the step that converts what was proved
into what was claimed, and omitting it leaves the proof formally incomplete. It
is also the step students most often skip, because by that point the arithmetic
feels finished. It isn't: the arithmetic established $\lnot Q\rightarrow\lnot
P$, and only the truth table in \emph{Why there are three methods at all} turns
that into $P\rightarrow Q$.
\end{remark}

% ===========================================================================
\method{3}{contradiction}{Assume $P$ and $\lnot Q$; produce a contradiction.}

\begin{nts}
Assuming both that $n$ is even and that $n+3$ is not odd, derive a statement of
the form $R$ and $\lnot R$. Nothing specifies which $R$ --- any contradiction
will do, and finding one is the creative part. Here the target will be a
violation of Fact~2: an integer exhibited as both even and odd.
\end{nts}

\begin{proof}[Proof 3]
Let $n\in\mathbb{Z}$ be arbitrary.
\begin{pstep}
Suppose $n$ is even and $n+3$ is not odd.
\begin{pstep}
$n=2k$ for some $k\in\mathbb{Z}$. \by{definition of even}

$n+3$ is even. \by{Fact 1: not odd, so even}

$n+3=2m$ for some $m\in\mathbb{Z}$. \by{definition of even}

$2k+3=2m$. \by{substituting $n=2k$}

$3=2m-2k=2(m-k)$. \by{algebra}

$m-k\in\mathbb{Z}$, so $3$ is even. \by{closure; definition of even}

$3=2\cdot 1+1$ and $1\in\mathbb{Z}$, so $3$ is odd.
\by{definition of odd}

So $3$ is both even and odd, contradicting Fact~2. \contra
\end{pstep}
Thus $n$ even and $n+3$ not odd is impossible.
\end{pstep}
Thus $\lnot(P\land\lnot Q)$ holds for every $n\in\mathbb{Z}$; equivalently, for
every $n\in\mathbb{Z}$, if $n$ is even then $n+3$ is odd.
\end{proof}

\begin{remark}
Two hypotheses enter and both get used --- $n=2k$ and $n+3=2m$ --- and
subtracting one from the other is what makes $n$ disappear, leaving a false
statement about the constant $3$. That is characteristic: a contradiction proof
often ends far from where it started, on a claim about something the original
statement never mentioned.
\end{remark}

\section*{Comparison}

\begin{center}
\begin{tabular}{lccc}
\toprule
 & Assume & Produce & Background used\\
\midrule
Direct         & $P$             & $Q$             & ---\\
Contrapositive & $\lnot Q$       & $\lnot P$       & Facts 1 and 2\\
Contradiction  & $P\land\lnot Q$ & \contra         & Facts 1 and 2\\
\bottomrule
\end{tabular}
\end{center}

\noindent
Three observations worth making out loud.

\begin{enumerate}[label=(\arabic*),leftmargin=2.4em,topsep=.4em,itemsep=.5em]
  \item \emph{The proofs are not equally cheap.} The direct proof spends
        nothing but the definitions. The other two spend both facts, because
        each begins with a negative hypothesis and a negative hypothesis gives
        you nothing to compute with until it has been turned positive. On this
        claim, then, the direct proof is simply the right one; the others are
        run here for practice.

  \item \emph{Contradiction hands you the most and demands the least.} You may
        assume $P$ and $\lnot Q$ both, and any absurdity discharges the
        obligation. That generosity is why it is the natural default when the
        conclusion is itself a negative or an impossibility --- as in the
        lexicographic argument later in the course, where the conclusion is that
        no utility representation exists and there is nothing to construct.
        The same generosity is why it is overused: an argument that could have
        been direct is often written by contradiction out of habit, with the
        second hypothesis never touched. If you finish a proof by contradiction
        without having used $\lnot Q$, you have written a direct proof inside a
        wrapper --- delete the wrapper.

  \item \emph{The closing line is part of the proof.} Methods 2 and 3 each end
        by converting what was established into what was claimed. That
        conversion is licensed by the truth table, not by the algebra, and a
        write-up that stops at the algebra has proved the wrong statement.
\end{enumerate}

\begin{remark}
A useful exercise, and the natural next one: run the same three methods on
\emph{if $n^{2}$ is odd, then $n$ is odd}. There the asymmetry is much sharper
--- the contrapositive proof is three lines, and the direct proof is genuinely
awkward. Comparing the two examples is what makes the point that choosing a
method is a real decision rather than a formality.
\end{remark}

\end{document}
